\newcommand{\DoNotLoadEpstopdf}{}
\documentclass[11pt,reqno]{amsart}
\usepackage[T1]{fontenc}
\usepackage{lmodern,amsmath,amssymb,amsthm,mathtools,mathrsfs,microtype}
\usepackage[margin=1.1in]{geometry}
\usepackage[hidelinks]{hyperref}
\hypersetup{pdftitle={Collision records, arithmetic, and raw local inversion in a triangular Lorentz gas},pdfauthor={Qian Qi}}
\setlength{\emergencystretch}{2em}
\allowdisplaybreaks[2]
\numberwithin{equation}{section}
\newtheorem{theorem}{Theorem}[section]
\newtheorem{lemma}[theorem]{Lemma}
\newtheorem{proposition}[theorem]{Proposition}
\newtheorem{corollary}[theorem]{Corollary}
\theoremstyle{definition}
\newtheorem{definition}[theorem]{Definition}
\newtheorem{inputstatement}[theorem]{Input statement}
\theoremstyle{remark}\newtheorem{remark}[theorem]{Remark}
\newcommand{\R}{\mathbb R}\newcommand{\Z}{\mathbb Z}
\newcommand{\one}{\mathbf1}\newcommand{\dd}{\,\mathrm d}
\newcommand{\esssup}{\operatorname*{ess\,sup}}
\newcommand{\diver}{\operatorname{div}}\newcommand{\TV}{\operatorname{TV}}
\title[Collision records and raw local inversion]{Collision records, arithmetic, and raw local inversion in a triangular Lorentz gas}
\author{Qian Qi}
\date{October 12, 2026. A2-DYN, revision 94}
\subjclass[2020]{37D50, 37A50, 60F05, 28A75}
\keywords{Lorentz gas, exact collision records, entrance traces, radial coarea, signed forcing, pointwise local limits}
\begin{document}
\raggedbottom
\begin{abstract}
We study the upper pointwise local limit problem for the exact collision
record of actual returns in a triangular Lorentz gas. The complete
noncritical entrance source has an essential-height estimate at every
fixed positive residence cutoff, with a leading constant independent
of both cutoffs. Its zero-residence contribution is identified with
the original weighted short-segment source. We also estimate all
critical histories which traverse a prescribed radial shell before
the target roof. Their entrance and growth are retained as a signed
packet; its actual target height is bounded without a residence or
incidence cutoff. Finally, projecting the complete signed physical
balance before taking positive parts gives a forcing term smaller
than the separate positive entrance and growth budgets. Its vanishing
is equivalent to the original pointwise assertion, in the prescribed
order of limits. The quantitative entrance and radial-transit estimates
do not establish vanishing of the remaining critical-target and
zero-residence concentration. The original pointwise theorem and all
exact physical data remain unchanged.
\end{abstract}
\maketitle
\input{journal-problem-v93}

\subsection{Entrance geometry and signed upper forcing}
The positive scalar remainder makes the upper sign the outstanding
part of \eqref{eq:v76-target}. The previous signed prefix balance
identified two sufficient positive mechanisms, entrance traces and
critical positive growth. Those positive pieces may overestimate the
projected source, since physical and prefix losses can cancel them.
The present argument retains that cancellation while estimating two
complete geometric families in essential height.

Theorem~\ref{thm:v94-entrance-height} controls every noncritical inlet
at each fixed positive residence cutoff by $C(r^{1/32}+r)$, uniformly
in that cutoff and the original guard. Proposition~\ref{prop:v94-inlet-charge}
identifies the precise uncut remainder as an original collision-source
restriction. Theorem~\ref{thm:v94-transit} controls the actual target
source of every history crossing the shell $\sqrt r\le z\le2\sqrt r$
within the backward window by
\[
 C r^{129/256}+C(1+\varepsilon^{-1})r^{257/256}.
\]
This includes genuine critical directional births; zero area at a
critical point is never used to erase its trace.

Theorems~\ref{thm:v94-signed-comparison} and \ref{thm:v94-endpoint}
then give a smaller, signed pointwise forcing and an equivalent upper
criterion. These are not finite-norm embeddings. The forcing remains
inside the essential supremum and collision limsup before its gradient
cutoff is removed. Its zero value is not proved by the two partial
height estimates.

\input{journal-inputs-v94}
\input{core/223_noncritical_entrance_residence}
\input{core/224_critical_radial_transit_packets}
\input{core/225_signed_projected_forcing_and_endpoint}

\section{The remaining upper-source estimate}
The entry theorem uses a complete physical-volume anchor for each
selected segment and only a finite residence-band sum before the
collision limit. The radial theorem obtains its anchor length from
actual passage through a gradient shell. Neither argument sums a
single-sheet modulus by an uncontrolled transverse multiplicity.
The signed theorem removes completed prefixes and preserves all
negative physical and artificial-cut contributions through projection.

The pointwise target nevertheless still requires vanishing of the
signed forcing in \eqref{eq:v94-equivalence}. Equation
\eqref{eq:v94-localization} records what is not estimated: a target
which remains in the small-gradient core, and the original weighted
zero-residence inlet concentration. Small fixed-window mass of these
sources does not bound their essential height. The equality of the
inlet and original-source charges is not a proof that either is zero.
Nor does a signed cancellation identity alone supply a bound for its
remaining positive part.

All original labels, section decisions, first-hit tests, weights and
arithmetic classes remain in \eqref{eq:v76-target}. The earlier lower
pointwise half, finite-norm laws and their conditioned consequences
retain their exact hypotheses. No point-conditioned likelihood, kernel,
unrestricted ratio or normalization of arithmetic zero classes is
promoted. The complete preceding article and every inherited proof
remain in the separate technical archive.

\appendix
\input{appendices/v94_dependency_map}
\input{journal-references-v93}
\end{document}
